% Generated draft: requires mathematical and source review.
\documentclass[11pt]{article}
\usepackage[margin=25mm]{geometry}
\usepackage{fontspec}
\setmainfont{Latin Modern Roman}
\usepackage{amsmath,amssymb,mathtools,unicode-math}
\setmathfont{Latin Modern Math}
\usepackage{enumitem,needspace,xurl,hyperref}
\setlength{\parindent}{0pt}
\setlength{\parskip}{5pt}
\setlength{\emergencystretch}{3em}
\newcommand{\N}{\mathbb{N}}
\newcommand{\Z}{\mathbb{Z}}
\newcommand{\Q}{\mathbb{Q}}
\newcommand{\R}{\mathbb{R}}
\newcommand{\C}{\mathbb{C}}
\newcommand{\F}{\mathbb{F}}
\newcommand{\A}{\mathbb{A}}
\newcommand{\PP}{\mathbb{P}}
\newcommand{\E}{\mathbb{E}}
\newcommand{\eps}{\varepsilon}
\newcommand{\dd}{\,\mathrm{d}}
\begin{document}
\section*{UnsolvedMath ID: 30004116}
\textbf{Source problem code:} catalogue entry
\subsection{Definitions and mathematical statement}
For an integer \(n\geq 4\), let
\[
q_n(x)=x_1^2+\cdots+x_n^2-x_{n+1}^2,
\qquad
\operatorname{PO}(n,1)=\operatorname{O}(q_n;\R)/\{\pm I\}.
\]
A lattice in this real Lie group is a discrete subgroup \(\Gamma\) for which \(\Gamma\backslash\operatorname{PO}(n,1)\) has finite invariant Haar volume. No torsion-free or cocompactness hypothesis is imposed in the supplied statement.

The question recorded in the source is:
\[
\text{For every integer }n\geq4\text{ and every lattice }
\Gamma\subset\operatorname{PO}(n,1),
\quad\text{is }\Gamma\text{ pseudo-arithmetic?}
\]
The supplied record does not define \emph{pseudo-arithmetic}. Its intended technical meaning must be recovered from the cited source before this question becomes a fully specified mathematical statement. The displayed formulation preserves the explicitly stated real-hyperbolic target; the accompanying literature assessment instead concerns complex-hyperbolic lattices. Those formulations cannot be identified from the supplied evidence. The record also does not specify an identity-component convention for the ambient group; the notation above uses the full projective orthogonal group.
\subsection{Short English statement}
The supplied question asks whether every finite-covolume discrete subgroup of the real projective orthogonal group \(\operatorname{PO}(n,1)\), in every integer dimension \(n\geq4\), is pseudo-arithmetic. The technical definition of pseudo-arithmeticity is missing, and the accompanying complex-hyperbolic literature assessment does not match the stated real-hyperbolic setting.
\subsection{Sources}
\begin{itemize}
\item Local UnsolvedMath registry: permanent ID \(30004116\), source problem code \texttt{catalogue entry}, title ``Pseudo-Arithmeticity of Higher-Dimensional Real-Hyperbolic Lattices,'' in the supplied Oberwolfach Reports Open Problems collection. The record gives proposed year 2019 and names no individual proposer.
\item Supplied report citation: \emph{Mini-Workshop: Reflection Groups in Negative Curvature} (2020), \emph{Oberwolfach Reports}, pp.~1043--1070, DOI 10.4171/owr/2019/17. \url{https://doi.org/10.4171/owr/2019/17}. The background describes Oberwolfach Reports as documenting mathematical workshops and their open-problem discussions. This supplied citation has not been independently inspected.
\item Supplied literature reference, labeled ``Pseudo-arithmetic complex hyperbolic lattices'': \url{https://arxiv.org/abs/1810.12837}. The supplied label and its relevance to the real-hyperbolic statement require verification.
\item Supplied literature reference, labeled ``Oberwolfach reflection-group report'': \url{https://doi.org/10.4171/owr/2019/17}.
\item Literature assessment dated \texttt{2026-08-22}: the supplied curation labels the problem open, asserts that known reflection and glued-lattice constructions are pseudo-arithmetic, and says that the cited literature leaves pseudo-arithmeticity of arbitrary complex-hyperbolic lattices unresolved. This is supplied dated curation, not an independently verified assessment. The background explicitly calls it literature triage derived from the supplied OWR curation and requests independent verification of time-sensitive status and cited claims.
\item The registry labels the statement \texttt{corrected\_verified} and says its mathematical content was preserved while notation, typography, duplicate clauses, or omitted definitions were repaired for self-contained publication. These are supplied editorial assertions; they do not resolve the missing definition or the mismatch of geometric settings.
\end{itemize}
\end{document}
